% Where You Adapt Beats What You Adapt
% Build:  pdflatex arc && pdflatex arc        (needs pgfplots >= 1.16)
% Figure data in figdata/*.dat, regenerated by: python3 paper/make_figures.py
% Nothing in this file is hand-copied from the lab notes; every plotted number comes
% from that script, which reads results/ directly.
\documentclass[11pt]{article}

\usepackage[margin=1in]{geometry}
\usepackage{booktabs}
\usepackage{amsmath}
\usepackage{graphicx}
\usepackage{pgfplots}
\usepackage[colorlinks=true,allcolors=blue]{hyperref}
\pgfplotsset{compat=1.16}
\usepackage{xcolor}

\definecolor{attncol}{HTML}{1B6CA8}
\definecolor{mlpcol}{HTML}{C8553D}
\definecolor{refcol}{HTML}{777777}

\newcommand{\dd}[2]{$#1 \pm #2$}
\newcommand{\TODO}[1]{\textcolor{red}{[\textbf{TODO:} #1]}}

\title{Where You Adapt Beats What You Adapt:\\
Depth Dominates Parameter Selection in Cross-Domain Few-Shot Adaptation}
\author{Ariet Jha}
\date{Draft --- \today}

\begin{document}
\maketitle

\begin{abstract}
Parameter-efficient adaptation is usually framed as a choice of \emph{which kind} of
parameter to update: normalisation affines, attention projections, adapters, low-rank
factors. We show that on a ViT-B/16 under cross-domain few-shot evaluation this framing
hides a larger effect. Holding parameter type and parameter count exactly fixed and
varying only \emph{which transformer block} is adapted, accuracy spans $8.9$ points on
EuroSAT --- against $6.5$ points for the entire parameter-type axis and $1.2$ points for
the attention-versus-MLP comparison inside it. Adapting attention in one well-chosen
block ($2.75\%$ of the backbone) beats full fine-tuning of all $85.8$M parameters by
\dd{4.43}{0.90} points; the same experiment in the block MLP ($5.50\%$) beats it by
\dd{5.72}{0.26}. The effect is not a property of attention --- the MLP depth curve has
the same shape, a comparable span and a higher peak --- and not a property of the
backbone alone: on ISIC, where the shift is larger, the curve is \emph{mirrored},
shallow blocks win, and no single block beats full fine-tuning at all. Four
shallow-versus-deep comparisons take opposite signs on the two datasets, each settled
across three episode draws.

We then ask why full fine-tuning loses, and test three pre-registered predictions of an
over-optimisation account. Two hold: reducing inner steps helps the large parameter
group and hurts the small one, and each method's gain over a frozen backbone grows with
support size at different rates --- full fine-tuning is \emph{worse than frozen} at
$k=1$ ($-1.20$), recovers to $+3.85$ at $k=5$ against attention's $+6.35$, and catches
up entirely by $k=20$ ($+7.41$ against $+7.65$). One fails: adding weight decay to the
inner loop does not repair full fine-tuning but degrades it further
(\dd{-1.08}{0.89}). We report the failure rather than reinterpret it, and note that the
advantage is therefore bounded to a \emph{regime} --- support sets too small to
constrain full fine-tuning but large enough for a small, well-placed group to learn
from --- which predicts the method's own obsolescence at larger $k$.

All numbers are validation-split; the test split is held for a single final run after
every design decision is frozen, and is deliberately absent from this paper.
\end{abstract}

\section{Introduction}

Cross-domain few-shot learning asks a pretrained backbone to classify from a handful of
labelled examples in a domain it was not trained on. The dominant design question in the
recent literature is \emph{which parameters to update}: LN-Tune updates LayerNorm
affines, SSF learns per-feature scale and shift, adapters and LoRA insert low-rank
structure. These are all answers to a question about parameter \emph{type}.

This paper reports that on a Vision Transformer the type question is secondary to a
question the type decomposition cannot express: \emph{at what depth}. A ViT-B/16 has
twelve blocks identical in shape and parameter count, so ``adapt block $i$ and nothing
else'' varies position with size held exactly fixed. A ResNet cannot support the same
experiment --- its stages differ in width, depth and resolution, so position is
confounded with size --- which is one reason the effect has gone unmeasured.

Our contributions:
\begin{enumerate}
  \item \textbf{Depth dominates type.} Five interventions identical in size and
    parameter type span $8.9$ points across depth on EuroSAT, against $6.5$ for the
    whole type axis (\S\ref{sec:depth}).
  \item \textbf{One well-placed block beats full fine-tuning} by $4.4$--$5.7$ points,
    and a badly placed block of the same size and type does not beat it at all. The win
    is a property of \emph{placement}, not of small interventions as such.
  \item \textbf{It is depth, not attention.} The same sweep in the block MLP reproduces
    the shape, the span and the interior peak (\S\ref{sec:mlp}), against a
    pre-registered prediction that it would not.
  \item \textbf{The curve mirrors with shift magnitude.} On ISIC the ordering reverses;
    four shallow-versus-deep comparisons have opposite signs on the two datasets
    (\S\ref{sec:mirror}).
  \item \textbf{A mechanism tested three ways, two confirmed and one failed.} Full
    fine-tuning over-optimises where there is little to gain; fewer inner steps and
    larger support sets both move it as predicted, while weight decay --- the
    intervention that should have repaired it --- makes it worse
    (\S\ref{sec:mech}). We report the failure as a failure.
  \item \textbf{A regime, not a universal win.} The single-block advantage requires a
    support set too small to constrain full fine-tuning and large enough for a small
    group to learn from. Outside that window the methods are indistinguishable: at
    $k=1$ nothing adapts usefully at all, and by $k=20$ full fine-tuning has caught up
    (\S\ref{sec:mech}). This bounds the claim and says when to use it.
\end{enumerate}

\section{Setup}

\paragraph{Protocol.} 5-way 5-shot, 15 queries per class, $100$ episodes per draw,
three draws (episode-sampling seeds $0,1,2$) at a fixed train/val/test partition.
Inner loop: $20$ Adam steps on the support set, then a cosine ProtoNet head~\cite{snell2017} on the
query set. Backbone ViT-B/16 \texttt{augreg\_in21k}~\cite{dosovitskiy2021,steiner2022}:
a standard supervised ImageNet-21k checkpoint, chosen for being strong and public. We
state plainly that it is \emph{our} choice and not a shared setting with any comparator
(Section~\ref{sec:related}).

\paragraph{Datasets and splits.} EuroSAT~\cite{helber2019} (overhead RGB, $10$ classes,
$27{,}000$ tiles) and ISIC~2018~\cite{tschandl2018,codella2019} (dermoscopy, $7$ classes,
$10{,}015$ images), both target sets of the BSCD-FSL cross-domain
benchmark~\cite{guo2020}. CropDiseases was out of scope and ChestX was unobtainable, so
the shift ladder has two rungs rather than four --- see Section~\ref{sec:limitations}.

Both are partitioned by \emph{group} at a fixed partition seed, at fractions chosen per
dataset rather than shared: EuroSAT $0.2/0.1/0.7$ train/val/test, giving
$5{,}400/2{,}700/18{,}900$ images, and ISIC $0.1/0.4/0.5$, giving
$1{,}000/4{,}013/5{,}002$ images over $747/2{,}988/3{,}735$ lesions. All classes are
present in every split of both. The fractions differ because ISIC's smallest classes are
small enough that a large training share leaves too few groups per class to sample
$5$-way episodes from the remainder; the asymmetry is deliberate and matches the
ResNet arm. Note that validation is the \emph{larger} partition on ISIC --- episodes are
drawn from $4{,}013$ images --- while on EuroSAT it is $2{,}700$.
The grouping key differs, and the difference matters. On ISIC it is the \texttt{lesion\_id}
from the HAM10000 metadata, so the $10{,}015$ images collapse to $7{,}470$ lesions
($4{,}482/1{,}494/1{,}494$) and group-disjointness is load-bearing: roughly a quarter of
images share a lesion with another image, and splitting by image would put two views of
one lesion on opposite sides of the partition. On EuroSAT each tile is an independent
sample, so the group is the image and the group-disjoint split \emph{reduces to} a plain
random image split. We note this rather than describing both as ``partitioned by
specimen'', which would be true of only one of them.

\paragraph{Statistics, and why they are stated explicitly.}
Our episode sampler seeds \emph{per episode}, so within a draw every baseline is
evaluated on literally the same episodes and all comparisons are \textbf{paired}.
Combining two within-run confidence intervals in quadrature --- which assumes
independence --- understated real margins by factors of $2.6$ and $4$ in our own early
analysis and produced one false null. We therefore report, for three-draw comparisons,
the per-seed difference as $\bar{d} \pm t_{2}\,s_d/\sqrt{3}$ with $t_2 = 4.303$.
Separately, repeated runs of byte-identical configurations differ by $\approx 0.6$
points (GPU floating-point non-determinism amplified through Adam), so we treat
$0.6$ as a floor and do not claim differences below it. That floor is entirely
attributable to the optimiser: a frozen backbone reproduces all $100$ episodes
bit-identically between runs, whereas two runs that take gradient steps agree exactly on
only $7$--$14$ of them while still averaging to within a point. The perturbation is
therefore near-independent across episodes, which is why per-episode pairing buys as
much as it does.

\section{Depth dominates parameter type}
\label{sec:depth}

Let \texttt{attn\_block}$i$ adapt the attention projections of block $i$ alone:
$2.75\%$ of the backbone, identical at every $i$ (asserted by test, since otherwise an
accuracy difference across depth could be about size). Figure~\ref{fig:depth} plots
accuracy against $i$.

\begin{figure}[t]
\centering
\begin{tikzpicture}
\begin{axis}[width=0.78\textwidth, height=6.2cm,
  xlabel={transformer block adapted}, ylabel={5-shot accuracy (\%)},
  xtick={0,2,5,8,11}, legend pos=south east, legend cell align=left,
  grid=major, grid style={gray!20}, ymin=76, ymax=91]
\addplot[attncol, mark=*, thick, error bars/.cd, y dir=both, y explicit]
  table[x=depth, y=euro_attn, y error=euro_attn_sd] {figdata/depth_curves.dat};
\addlegendentry{EuroSAT, attention ($2.75\%$)}
\addplot[mlpcol, mark=square*, thick, dashed, error bars/.cd, y dir=both, y explicit]
  table[x=depth, y=euro_mlp, y error=euro_mlp_sd] {figdata/depth_curves.dat};
\addlegendentry{EuroSAT, block MLP ($5.50\%$)}
\addplot[refcol, dotted, thick, domain=0:11, samples=2]{83.07};
\addlegendentry{full fine-tuning ($100\%$)}
\addplot[refcol, dashdotted, thick, domain=0:11, samples=2]{79.22};
\addlegendentry{frozen backbone}
\end{axis}
\end{tikzpicture}
\caption{EuroSAT depth curves for two parameter types, three draws, error bars
$\pm 1$ s.d.\ across draws. Both types rise to an interior peak at block 8 and fall at
block 11. A single block of either type beats full fine-tuning; a single block at the
shallow end does not. The argmax is block 8 in all three draws for both types.}
\label{fig:depth}
\end{figure}

\paragraph{The span.} Across the five depths the attention curve spans $8.9$ points
($8.89/9.31/8.41$ in the three draws individually). For comparison, the whole
parameter-type axis --- patch-embed, LayerNorm, attention, MLP, everything --- spans
$6.5$ points, and attention-versus-MLP inside it is \dd{1.18}{0.43}. Depth is the
larger axis.

\paragraph{The peak is interior.} Block 8 beats block 11 by \dd{1.54}{0.80}, every draw
positive, and beats block 5 by $5.71 \pm 1.38$. The shallow pair is
indistinguishable: block 2 $-$ block 0 is $-0.12 \pm 0.48$.

\paragraph{One block beats everything larger.}
Block 8 alone beats full fine-tuning by \dd{4.43}{0.90}, the frozen backbone by
\dd{8.28}{1.27}, and \emph{all twelve blocks} of attention by \dd{1.93}{0.63}.
Figure~\ref{fig:params} makes the point that matters for practice: the relationship
between adapted fraction and accuracy is not monotone, and full fine-tuning sits below
interventions $20$--$40\times$ smaller.

\begin{figure}[t]
\centering
\begin{tikzpicture}
\begin{axis}[width=0.78\textwidth, height=5.6cm,
  xlabel={parameters adapted (\% of backbone, log scale)},
  ylabel={5-shot accuracy (\%)}, xmode=log, grid=major, grid style={gray!20},
  nodes near coords={\footnotesize\pgfplotspointmeta}, point meta=explicit symbolic,
  every node near coord/.append style={anchor=south, yshift=2pt}]
\addplot[only marks, mark=*, attncol, mark size=2.2pt]
  table[x=frac, y=acc, meta=label] {figdata/params_vs_acc.dat};
\end{axis}
\end{tikzpicture}
\caption{Adapted fraction against accuracy on EuroSAT (log $x$; the frozen point at
$0\%$ is omitted from the axis and annotated in text). More is not better: the best
configuration adapts $5.5\%$ of the backbone.}
\label{fig:params}
\end{figure}

\section{It is depth, not attention}
\label{sec:mlp}

The result above could be a fact about attention rather than about depth. We tested this
by running the identical sweep in the block MLP, and recorded both predictions in the
code before the measurement existed: if the effect is \emph{about attention}, MLP should
stay near its whole-type value at every depth, with no comparable spread and no single
MLP block beating full fine-tuning by anything like $5$ points; if it is \emph{about
depth}, MLP should show the same shape.

The second branch is what happened. MLP in block 8 scores the highest accuracy we
measure anywhere and beats full fine-tuning by \dd{5.72}{0.26}; the MLP curve spans
$8.5$ points against attention's $8.9$; its argmax is block 8 in all three draws.
Depth is about as decisive inside the MLP as inside attention, in a parameter type
twice the size.

\paragraph{Type still matters, as a magnitude.} MLP is ahead of attention at every
depth measured and settled-behind at none, by a margin running \dd{1.29}{0.68} at
block 8 to $+4.66 \pm 1.84$ at block 5 --- a factor of $3.6$. An earlier single-draw
reading had this margin changing \emph{sign} at block 11; three draws demote that to
$-0.96 \pm 1.61$, unsettled, and we do not claim it.

\paragraph{A corollary about whole-type comparisons.} Whole-type attention beats
whole-type MLP by \dd{1.18}{0.43}, yet attention never demonstrably beats MLP at any
single depth. Whole-type MLP adapts $66\%$ of the backbone against attention's $33\%$
and suffers more from adapting twelve blocks at once; that penalty is what the
whole-type average measures. \emph{Attention-versus-MLP is a fact about adapting all
twelve blocks simultaneously and carries no information about which type is better to
adapt at a given depth.}

\section{The curve mirrors with shift magnitude}
\label{sec:mirror}

On ISIC the same five depths give the opposite ordering: the shallow half is best and
the deep end is worst (Figure~\ref{fig:mirror}).

\begin{figure}[t]
\centering
\begin{tikzpicture}
\begin{axis}[width=0.78\textwidth, height=5.4cm,
  xlabel={transformer block adapted}, ylabel={5-shot accuracy (\%), ISIC},
  xtick={0,2,5,8,11}, grid=major, grid style={gray!20}]
\addplot[attncol, mark=*, thick, error bars/.cd, y dir=both, y explicit]
  table[x=depth, y=isic_attn, y error=isic_attn_sd] {figdata/depth_curves.dat};
\end{axis}
\end{tikzpicture}
\caption{ISIC depth curve, attention, three draws. A flat plateau over blocks $0$--$5$
(no internal comparison settled, and the argmax moves between draws) then a step down
to the deep end. Compare Figure~\ref{fig:depth}: the ordering is reversed.}
\label{fig:mirror}
\end{figure}

\begin{table}[t]
\centering
\begin{tabular}{lrr}
\toprule
shallow $-$ deep & ISIC & EuroSAT \\
\midrule
block 2 $-$ block 11 & \dd{3.47}{1.14} & \dd{-7.31}{1.13} \\
block 5 $-$ block 11 & \dd{3.28}{2.07} & \dd{-4.17}{1.31} \\
block 5 $-$ block 8  & \dd{1.93}{0.23} & \dd{-5.71}{1.38} \\
\bottomrule
\end{tabular}
\caption{Every shallow-versus-deep comparison has the opposite sign on the two
datasets, all settled, three draws per side, same selector and protocol throughout.}
\label{tab:mirror}
\end{table}

Two differences from EuroSAT are worth stating. ISIC's depth span is $3.7$ points
against EuroSAT's $8.9$ --- real but less than half the size --- and \emph{no} single
block beats full fine-tuning there (the best loses by \dd{3.72}{1.03}). So the
\emph{depth effect} is a property of the backbone, while the \emph{single-block win} is
specific to mild shift.

\paragraph{Interpretation.} We read the depth curve as locating where the pretrained
representation stops transferring. Dermoscopy diverges from ImageNet-21k early, so ISIC
needs early adaptation; EuroSAT is overhead RGB, where low-level features carry over and
only the task-specific top needs moving. Independent support, not from accuracy: the
same selector at the same depth tunes to opposite ends of the same learning-rate grid on
the two datasets --- block 0 wants $3\times10^{-5}$ on EuroSAT and $3\times10^{-3}$ on
ISIC.

\section{Mechanism}
\label{sec:mech}

\S\ref{sec:depth} leaves open why full fine-tuning loses at all. We propose that
$85.8$M parameters taking $20$ steps on $25$ support images over-optimise where there is
little to gain, and that the penalty scales inversely with how much there is to gain.
This predicts specific, checkable things.

\paragraph{Confirmed: less optimisation helps the larger group.}
Figure~\ref{fig:steps} shows accuracy against inner steps. Attention improves
monotonically; full fine-tuning is $1.40$ points \emph{worse} at $20$ steps than at
$10$. Best-against-best, attention still wins by $1.32$.

\begin{figure}[t]
\centering
\begin{tikzpicture}
\begin{axis}[width=0.7\textwidth, height=5cm,
  xlabel={inner adaptation steps}, ylabel={5-shot accuracy (\%)},
  xtick={5,10,20}, grid=major, grid style={gray!20}, legend pos=south west]
\addplot[attncol, mark=*, thick] table[x=steps, y=attn] {figdata/inner_steps.dat};
\addlegendentry{attention only ($33\%$)}
\addplot[refcol, mark=square*, thick, dashed] table[x=steps, y=full] {figdata/inner_steps.dat};
\addlegendentry{full fine-tuning ($100\%$)}
\end{axis}
\end{tikzpicture}
\caption{More optimisation helps the small group and hurts the large one. EuroSAT,
seed $0$ at all three step counts.}
\label{fig:steps}
\end{figure}

\paragraph{Confirmed: the deficit shrinks as the support set grows.}
The account predicts that full fine-tuning's penalty is worst where there is least to
fit, so its \emph{gain over a frozen backbone} should grow with $k$. It does, and the
shape is informative (Figure~\ref{fig:shots}, Table~\ref{tab:gains}). At $k=1$ full
fine-tuning is \emph{below} frozen; attention is indistinguishable from frozen, because
one example per class supports no learning for either. At $k=5$ attention gains
\dd{6.35}{1.07} against full fine-tuning's \dd{3.85}{0.43}, and the \dd{2.50}{0.64}
margin is that difference. At $k=20$ attention has saturated (\dd{7.65}{0.98}) while full
fine-tuning has caught up (\dd{7.41}{1.04}).

\begin{table}[t]
\centering
\begin{tabular}{lrrr}
\toprule
gain over frozen & $k=1$ & $k=5$ & $k=20$ \\
\midrule
attention only ($33\%$)  & \dd{0.05}{1.42} & \dd{6.35}{1.07} & \dd{7.65}{0.98} \\
full fine-tuning ($100\%$) & \dd{-1.20}{1.42} & \dd{3.85}{0.43} & \dd{7.41}{1.04} \\
\midrule
attention $-$ full        & \dd{1.25}{1.18} & \dd{2.50}{0.64} & \dd{0.24}{0.53} \\
\bottomrule
\end{tabular}
\caption{EuroSAT. The $k=5$ row is three draws; $k=1$ and $k=20$ are one draw each and
are paired per episode. Only the $k=5$ difference and the two $k=20$ gains clear both
their interval and the $0.6$ floor.}
\label{tab:gains}
\end{table}

\begin{figure}[t]
\centering
\begin{tikzpicture}
\begin{axis}[width=0.72\textwidth, height=5.4cm,
  xlabel={support examples per class $k$ (log scale)}, xmode=log, log basis x=10,
  ylabel={5-way accuracy (\%)}, xtick={1,5,20}, xticklabels={1,5,20},
  grid=major, grid style={gray!20}, legend pos=south east, legend cell align=left]
\addplot[attncol, mark=*, thick] table[x=k, y=attn] {figdata/shots.dat};
\addlegendentry{attention only ($33\%$)}
\addplot[refcol, mark=square*, thick, dashed] table[x=k, y=full] {figdata/shots.dat};
\addlegendentry{full fine-tuning ($100\%$)}
\addplot[refcol, mark=triangle*, thick, dotted] table[x=k, y=frozen] {figdata/shots.dat};
\addlegendentry{frozen backbone}
\end{axis}
\end{tikzpicture}
\caption{The advantage is a window. All three curves converge at $k=1$, where nothing
adapts usefully, and the two adapted curves converge at $k=20$, where full fine-tuning
has enough support to constrain its parameter count. The gap closes because full
fine-tuning catches up, not because attention degrades --- attention never gets worse
at any $k$.}
\label{fig:shots}
\end{figure}

Note that the \emph{raw gap} is non-monotone in $k$ ($+1.25$, $+2.50$, $+0.24$) and
cannot be used to test this account at all: it is a difference of two gains that are
both floored at $k=1$ and both large at $k=20$, so it must fall at either end whatever
the mechanism. The gain over frozen is the statistic the account predicts, and we report
it in place of the gap for that reason.

\paragraph{Failed: weight decay does not repair full fine-tuning.}
The account's most direct prediction is that constraining the large group should help
it. Adding weight decay to the inner loop does the opposite, and does it as a
dose--response. Against its own zero-decay arm at the same tuned learning rate, full
fine-tuning loses \dd{-0.31}{0.89} at $10^{-4}$, \dd{-1.08}{0.89} at $10^{-3}$,
\dd{-3.44}{0.98} at $10^{-2}$ and \dd{-5.12}{1.07} at $10^{-1}$ --- strictly
decreasing, the last three settled. The prediction fails in direction, and the monotone
decline makes the failure decisive rather than marginal.

Attention is unaffected up to $10^{-2}$ (\dd{0.61}{0.99} at $10^{-3}$, \dd{-0.44}{1.03}
at $10^{-2}$, both unsettled) but is destroyed at $10^{-1}$ (\dd{-5.59}{1.19}), by as
much as full fine-tuning is. So the asymmetry is in \emph{where the damage threshold
sits}, not in whether a small group can be harmed by regularisation at all.

We record this as a failure rather than reinterpret it. There is a principled reason the
test may be the wrong one: weight decay pulls parameters toward \emph{zero}, and
pretrained weights are nowhere near zero, so in a fine-tuning inner loop
decay-toward-the-origin does not express ``stay near the prior'' but actively destroys
the features the checkpoint exists to provide. The regulariser matching the account's
intent is $\ell^2$ toward the pretrained weights, which we have not run. The
dose--response is consistent with that reading --- damage proportional to decay
strength is what dragging pretrained features toward zero would produce --- but the
explanation arrived after the result, and an explanation produced after a failure is a
hypothesis, not a finding. It will take the $\ell^2$-SP experiment to change this entry.

We also flag a confound in the other half. The inner optimiser applies \emph{coupled}
$\ell^2$ (decay added to the gradient, then normalised by Adam's second moment), and
the two arms run at tuned learning rates a factor of ten apart. Equal nominal decay is
therefore not equal regularisation strength across baselines, so ``decay hurts full
fine-tuning but not attention'' is not a controlled comparison and we do not claim
attention is insensitive to regularisation. The within-\texttt{full\_ft}
dose--response is clean, because nothing varies along it but the decay.

\paragraph{Scoreboard, and what the mechanism is worth.}
Two of three pre-registered predictions hold: inner steps and support size. One fails:
regularisation. The honest summary is that full fine-tuning behaves like a method with
too much capacity for its data --- it degrades with more optimisation, it is worse than
doing nothing at $k=1$, and it recovers as data arrives --- but the one intervention
that should have fixed it instead made it worse, and we cannot yet say whether that
indicts the account or the intervention.

The depth result does not rest on any of this. It requires full fine-tuning to lose at
$k=5$, which it does across three draws on both datasets, and not an explanation of
why.

\section{Relation to the parameter-type literature}
\label{sec:related}

On ResNet-50 we find what the type literature implies: a $0.75$-point band across a
$778\times$ range of trainable parameter counts, replicated on both datasets, with
adaptation $\approx 90\%$ a per-channel affine and near rank-$1$ at the top of the
network. We now read that null as \emph{regime-specific} rather than general: it was
measured on ISIC, the strong-shift setting, where our ViT results also show placement
mattering least and joint adaptation paying most.
\paragraph{We do not report comparator numbers, and why.}
The three methods this work is positioned against do not share a setting with us, and we
checked rather than assumed. SSF~\cite{lian2022} tunes per-channel scale and shift on a
ViT-B/16 ImageNet-21k checkpoint --- the same checkpoint family we use --- but evaluates
on FGVC and VTAB-1k, which are not episodic few-shot benchmarks. LN-Tune~\cite{basu2024}
tunes LayerNorm only, the direct analogue of the baseline our ResNet result turns on, and
reports on Meta-Dataset and ORBIT. DIPA~\cite{perera2024} uses ViT-S/16 with DINO
self-supervised pretraining: a different capacity \emph{and} a different pretraining
objective.

Backbone, pretraining objective, benchmark and episode protocol therefore all differ,
and a table placing their published accuracies beside ours would compare four things at
once while appearing to compare one. Re-running all three under our protocol is the
right experiment and we have not done it; until then we make no quantitative claim
against them. What we do claim is independent of them: it is an \emph{internal}
comparison, every arm of which is run by us under one protocol, in which placement
beats full fine-tuning and beats parameter-type choice. Our baseline for ``what the type
literature implies'' is our own ResNet replication, not their numbers.

\section{Limitations}
\label{sec:limitations}

\begin{itemize}
  \item \textbf{Validation split only.} Every number here is validation. The test-split
    run is deliberately held until all design decisions are frozen.
  \item \textbf{Five of twelve depths.} The curve between measured points is
    interpolation; whether the peak is at block 8 specifically or anywhere in $6$--$10$
    is not established.
  \item \textbf{Two datasets.} Two points cannot establish a monotone relationship
    between shift magnitude and anything. Headroom is our proxy for shift magnitude and
    is a proxy, not a measurement.
  \item \textbf{The type\,$\times$\,depth interaction is one dataset} --- it has not
    been measured on ISIC.
  \item \textbf{Runs are not reproducible to better than $\approx 0.6$ points}, and we
    have not enabled deterministic kernels. Every individual accuracy carries that noise
    on top of the episode-sample spread.
  \item \textbf{One architecture family} for the depth result.
  \item \textbf{The support-size result is EuroSAT only.} ISIC already inverts the
    depth ordering, so its support-size behaviour cannot be assumed to match; it is
    unmeasured. Within EuroSAT, $k=1$ and $k=20$ are a single draw each.
  \item \textbf{The gains-over-frozen statistic was chosen after seeing $k=1$ and
    $k=5$.} Only the $k=20$ row was out-of-sample for it. That is one genuinely
    predictive point, not a replicated test.
  \item \textbf{The failed regularisation prediction is unresolved.} We do not know
    whether weight decay's failure indicts the over-optimisation account or only the
    choice of regulariser, and the experiment that would separate them ($\ell^2$
    toward the pretrained weights) has not been run.
  \item \textbf{No comparator is re-run under our protocol.} Our positioning rests on
    an internal comparison and on our own ResNet replication, not on published numbers
    measured in other settings (\S\ref{sec:related}).
\end{itemize}

\section{Conclusion}

The parameter-efficient adaptation literature asks which kind of parameter to update.
On a ViT-B/16 under cross-domain few-shot evaluation, that question is dominated by one
it cannot express: at which depth. Five interventions identical in type and size span
$8.9$ points across depth, against $6.5$ for the whole type axis, and one well-placed
block of either attention or MLP beats full fine-tuning of the entire backbone while a
badly placed block of the same size and type does not beat it at all. The win belongs to
placement, not to small interventions as such.

Two features of the result bound it, and we think both are more useful than a larger
headline would have been. The depth curve \emph{reverses} between our two datasets, so
the right depth is a property of the domain shift rather than of the architecture --- we
read the curve as locating where the pretrained representation stops transferring, and
on the stronger-shift dataset no single block suffices. And the advantage occupies a
window in support size: it needs data too scarce to constrain full fine-tuning and
plentiful enough for a small group to learn from, so it disappears at both ends. A
method whose own obsolescence at $k=20$ is visible in its figures is easier to deploy
correctly than one whose limits are unstated.

What we cannot yet explain is why the obvious repair for full fine-tuning does not work.
Weight decay degrades it rather than rescuing it, which is either a fact about the
account or a fact about decaying toward the origin from a pretrained initialisation. The
experiment that distinguishes those is the natural next one, along with a second
architecture family and a third dataset to turn the mirror from an instance into a
trend.

\section*{Reproducibility}

Every figure in this paper is generated by \texttt{paper/make\_figures.py} directly from
the raw per-episode result files, with the number of draws emitted alongside each value.
No number in a figure or table was transcribed by hand. The lab record, including seven
instances of a single-draw result moving under additional draws (four promoted, three
demoted) and four retracted readings of our own data with the cause of each, is kept in
\texttt{docs/RESULTS\_VIT.md}.

\begin{thebibliography}{9}

\bibitem{basu2024}
S.~Basu, D.~Massiceti, S.~X.~Hu, and S.~Feizi.
\newblock Strong baselines for parameter-efficient few-shot fine-tuning.
\newblock In \emph{AAAI Conference on Artificial Intelligence}, 2024.
  arXiv:2304.01917.

\bibitem{codella2019}
N.~Codella et~al.
\newblock Skin lesion analysis toward melanoma detection 2018: A challenge hosted by the
  International Skin Imaging Collaboration (ISIC).
\newblock arXiv:1902.03368, 2019.

\bibitem{dosovitskiy2021}
A.~Dosovitskiy et~al.
\newblock An image is worth 16x16 words: Transformers for image recognition at scale.
\newblock In \emph{International Conference on Learning Representations}, 2021.

\bibitem{guo2020}
Y.~Guo et~al.
\newblock A broader study of cross-domain few-shot learning.
\newblock In \emph{European Conference on Computer Vision}, 2020.

\bibitem{helber2019}
P.~Helber, B.~Bischke, A.~Dengel, and D.~Borth.
\newblock EuroSAT: A novel dataset and deep learning benchmark for land use and land
  cover classification.
\newblock \emph{IEEE Journal of Selected Topics in Applied Earth Observations and Remote
  Sensing}, 2019.

\bibitem{lian2022}
D.~Lian, D.~Zhou, J.~Feng, and X.~Wang.
\newblock Scaling \& shifting your features: A new baseline for efficient model tuning.
\newblock In \emph{Advances in Neural Information Processing Systems}, 2022.
  arXiv:2210.08823.

\bibitem{perera2024}
R.~Perera and S.~Halgamuge.
\newblock Discriminative sample-guided and parameter-efficient feature space adaptation
  for cross-domain few-shot learning.
\newblock In \emph{IEEE/CVF Conference on Computer Vision and Pattern Recognition}, 2024.

\bibitem{snell2017}
J.~Snell, K.~Swersky, and R.~Zemel.
\newblock Prototypical networks for few-shot learning.
\newblock In \emph{Advances in Neural Information Processing Systems}, 2017.

\bibitem{steiner2022}
A.~Steiner, A.~Kolesnikov, X.~Zhai, R.~Wightman, J.~Uszkoreit, and L.~Beyer.
\newblock How to train your ViT? Data, augmentation, and regularization in vision
  transformers.
\newblock \emph{Transactions on Machine Learning Research}, 2022.

\bibitem{tschandl2018}
P.~Tschandl, C.~Rosendahl, and H.~Kittler.
\newblock The HAM10000 dataset, a large collection of multi-source dermatoscopic images
  of common pigmented skin lesions.
\newblock \emph{Scientific Data}, 5:180161, 2018.

\end{thebibliography}

\end{document}
